\section{TRABALHOS CORRELATOS}

Nesta seção, realiza-se uma análise breve de jogos sérios para o ensino de Química. A investigação concentra-se em objetivos pedagógicos, \textit{game design} e resultados alcançados, de modo a estabelecer as bases para a diferenciação técnica do projeto Alchemon.

\subsection{ChemCaper: Act I - Petticles in Peril}

O \textit{ChemCaper} é um jogo de \textit{Role-Playing Game} (RPG) de aventura, desenvolvido pela \textit{ACE EdVenture Studio} em colaboração com a \textit{Artoncode}. O objetivo central desse jogo é auxiliar estudantes na compreensão de conceitos do currículo \textit{International General Certificate of Secondary Education} (IGCSE) de Química, abrangendo o uso de aparelhos laboratoriais, técnicas de separação, periodicidade e ligações químicas.

Disponível na plataforma Steam, o jogo utiliza mecanismos de aprendizagem implícita e uma direção de arte atrativa para engajar o público jovem. Os resultados de estudos qualitativos indicam que o \textit{ChemCaper} funciona como um excelente gancho motivacional \cite{Lackey2022}. Estudantes que utilizaram o jogo demonstraram maior disposição para enfrentar problemas complexos de ligações químicas, os quais tradicionalmente geram desinteresse em formatos expositivos \cite{Lackey2022}. Além disso, o jogo foi capaz de modelar e prover \textit{cantilevers} para normas de laboratório que muitas vezes não podem ser executadas fisicamente em escolas com recursos limitados.

\subsection{ChemCraft}

O \textit{ChemCraft} é um jogo \textit{Creature Capture} 2D focado no ensino de Química Orgânica de nível superior (\textit{IB Higher Level}). O objetivo do projeto, desenvolvido por Jihae Han \cite{Han2021}, foi investigar a transposição de sistemas químicos reais para regras de sistemas digitais, seguindo o conceito de \textit{gaming literacy}.

As mecânicas de \textit{gameplay} do \textit{ChemCraft} seguem reações químicas reais, como combustão e síntese, nas quais as propriedades dos objetos do jogo referenciam valores de dados químicos reais, como massa atômica e energia de dissociação \cite{Han2021}. Os resultados de um estudo piloto revelaram que o \textit{software} é considerado útil para a compreensão de estruturas moleculares complexas, embora os autores tenham identificado que a progressão acelerada do conteúdo pode sobrecarregar a carga cognitiva do estudante iniciante \cite{Han2021}.

\subsection{Batalha QuímiCard}

No cenário educacional brasileiro, o Batalha QuímiCard é um jogo sério de cartas digital desenvolvido para a plataforma Android. O objetivo do jogo é validar o uso de mecânicas inspiradas em títulos comerciais para o ensino de ácidos e bases, concentrando-se especificamente nas forças relativas das substâncias e nas constantes de equilíbrio $K_a$ e $K_b$ \cite{GomesdeFreitas2017}.

A proposta fundamenta-se na aprendizagem tangencial \cite{Portnow2012}, processo no qual o estudante internaliza o conhecimento científico para obter vantagens estratégicas em duelos \cite{LimadaCostaJunior2022}. Os resultados de validação apontaram que 64\% dos estudantes consideraram o jogo útil para a compreensão de forças ácido-base, enquanto 40\% atingiram proficiência total na identificação de conceitos químicos fortes após a interação lúdica \cite{GomesdeFreitas2017}. O sucesso do projeto reforça a existência de uma demanda nacional por ferramentas digitais que transcendam o modelo tradicional de memorização mecânica.

\subsection{Inferência sobre os trabalhos correlatos}

A análise dos similares evidencia que o uso de mecânicas de RPG e colecionismo constitui uma estratégia para lidar com a abstração e complexidade da Química. Projetos como \textit{ChemCaper} e Batalha QuímiCard \cite{GomesdeFreitas2017} validam a viabilidade técnica de utilizar batalhas e exploração como motores de engajamento, ao passo que o \textit{ChemCraft} \cite{Han2021} demonstra a importância da fidelidade paramétrica aos dados científicos reais e uso de jogo do gênero \textit{Creature Capture}.
